\subsection{Shuffle matrices}
For each column $j$, choose uniformly an assignment of the symbols
$[k]_0$ to the $N$ row labels, with exactly $n_i^{(j)}$ occurrences of
symbol $i$. Choose these assignments independently across columns.
The resulting $N\times K$ matrix is uniform among all matrices with
the prescribed column counts. Write its row words as $w_1,\ldots,w_N$.
Stable sorting orders these words lexicographically and orders equal
words by their original row labels. Its permutation maps sorted
positions to original labels. Write $s_1\le_{\rm lex}\cdots\le_{\rm lex}s_N$
for the sorted words, and join $i$ to $i+1$ when $s_i=s_{i+1}$.
This defines the shuffle graph $G$.

Let $P_N^K$ denote the law of the inverse of the chronological
composition of the $K$ forward shuffles, started from the identity
deck. Thus the one-step inverse operations are applied in reverse
chronological order. Write $\nuUnif$ for uniform measure on
$\mathfrak S_N$, and
\[
 \dTV(P,Q)=\frac12\sum_{\sigma\in\mathfrak S_N}|P(\sigma)-Q(\sigma)|.
\]
For $H\subseteq\mathbb Z$, let $E(H)$ be its internal edges in the
integer path and let $\partial H$ consist of the edges with exactly
one endpoint in $H$. The latter includes edges at the ends of the
deck. For a graph $G$, $E(G)$ denotes its edge set; we also identify
$G$ with this set when taking intersections. For a permutation $\pi$,
write $\pi_G$ for the permutation obtained by sorting its values
increasingly inside every component of $G$.

\begin{lemma}[Matrix and graph representations]\label{lem:matrix-graph}
Stable sorting the fixed-count matrix has law $P_N^K$. If $G$ is
sampled from this matrix and $\pi$ is a uniform permutation independent
of the matrix, then $\pi_G$ also has law $P_N^K$. More generally,
for every real function $F$ of the sorted words and a permutation,
\[
 \E F\bigl((s_i)_{i=1}^N,\operatorname{sort}(w_1,\ldots,w_N)\bigr)
 =\E\E_\pi F\bigl((s_i)_{i=1}^N,\pi_G\bigr).
\]
The total-variation distances of the forward and inverse shuffle
laws from $\nuUnif$ are equal.
\end{lemma}
\begin{proof}
A one-step interleaving is specified by the word recording which
pile supplies each successive card. These words are uniform with
their prescribed counts. Its inverse gathers equal symbols in
increasing symbol order and preserves their previous order, which
is exactly stable sorting by that symbol.

For several steps, the symbols are initially assigned to the current
positions at each inverse operation. Process the inverse operations
from column $K$ back to column $1$, pulling each symbol assignment
back to the original row labels by the permutation already determined
by the later columns. This reindexes only that column and preserves
its counts. The reindexing permutation depends only on later columns,
so it is recovered when these later columns are known. The operation
therefore defines a bijection of the finite product of column-assignment
spaces, with inverse obtained by reversing each reindexing. It
preserves the uniform law on that product.

To identify the resulting permutation, use induction on the list of
columns. Sorting by a leading symbol, while preserving the order
already determined by the suffix, is the same as sorting by the
whole word. When two whole words agree, both procedures preserve
their original label order. The empty list gives the identity.
Consequently the pullback bijection takes the physical inverse
permutation to stable sorting of the full matrix. This proves the
first assertion, including the chronological convention.

Next relabel all rows by an independent uniform permutation. Every
fixed relabeling preserves the column counts and the uniform matrix
law, and it leaves the sorted words unchanged. Conditional on the
matrix, the relabeled original labels appear in a uniform order in
the sorted positions. Stable sorting restores increasing label
order precisely inside the equal-word components of $G$. Thus the
stable sort of the relabeled matrix is obtained by taking a uniform
permutation and sorting it inside those components. Averaging first
over the relabeling and then over the matrix proves the displayed
identity and the graph-mixture assertion.

Finally inversion is a bijection of $\mathfrak S_N$ preserving
uniform measure. Reindexing the finite sum defining total variation
by this bijection proves the last assertion.
\end{proof}
